% OSNIP-style task metrics: paper vs fastProve measured numbers
% Compile: xelatex osnip_vs_fastprove_tasks.tex
\documentclass[11pt,a4paper]{article}
\usepackage[margin=1.6cm]{geometry}
\usepackage{fontspec}
\usepackage{xeCJK}
\setmainfont{Times New Roman}
\setCJKmainfont{PingFang SC}
\usepackage{booktabs,array,xcolor,colortbl,amsmath,multirow}
\usepackage{hyperref}

\definecolor{headerbg}{HTML}{F0F0F0}
\definecolor{note}{HTML}{555555}

\title{\textbf{OSNIP 同协议任务精度：论文 vs fastProve}\\
\large Llama-3.2-3B-Instruct}
\author{fastProve evaluation}
\date{\today}

\begin{document}
\maketitle
\vspace{-1.0em}

\noindent
{\small\color{note}
\textbf{说明：}
左列取自 OSNIP 论文 Table~1 / Table~2（Llama-3.2-3B-Instruct）；
右列为本仓库当前协变混淆方案在 RTX~5090 上实测
（\texttt{results/raw/osnip\_style\_bench\_n200.json}，每任务 $n{=}200$，1 master key）。
评分：闭集任务为 continuation 对数似然 MC；WikiText-2 为滑动窗口 teacher-forced PPL。
\textbf{绝对分}受 prompt / chat 模板 / 子集影响，不可与论文硬对齐；
\textbf{相对保留率 RP} 与 \textbf{PPL 相对增幅} 才是公平主指标。
}

\vspace{0.7em}

\subsection*{闭集准确率与 RP}

\begin{center}
\renewcommand{\arraystretch}{1.18}
\footnotesize
\begin{tabular}{@{}lcccccc@{}}
\toprule
\rowcolor{headerbg}
任务 &
OSNIP 明文 & OSNIP & OSNIP RP &
fastProve 明文 & structural & full \\
\midrule
ARC-E
  & 0.712 & 0.707 & 99.3\%
  & 0.535 & 0.530 & 0.505 \\
PIQA
  & 0.768 & 0.766 & 99.7\%
  & 0.460 & 0.460 & 0.460 \\
MNLI
  & 0.542 & 0.535 & 98.7\%
  & 0.480 & 0.485 & 0.495 \\
SST-2
  & 0.867 & 0.825 & 95.2\%
  & 0.970 & 0.965 & 0.950 \\
ANLI-R1
  & 0.440 & 0.431 & 98.0\%
  & 0.475 & 0.480 & 0.415 \\
ANLI-R2
  & 0.401 & 0.407 & 101.5\%
  & 0.395 & 0.400 & 0.355 \\
ANLI-R3
  & 0.425 & 0.410 & 96.5\%
  & 0.350 & 0.350 & 0.350 \\
WiC
  & 0.497 & 0.513 & 103.2\%
  & 0.505 & 0.505 & 0.505 \\
HellaSwag
  & 0.716 & 0.709 & 99.0\%
  & 0.465 & 0.465 & 0.455 \\
MMLU
  & 0.606 & 0.572 & 94.4\%
  & 0.295 & 0.295 & 0.315 \\
\midrule
\textbf{闭集平均}
  & 0.597 & 0.588 & \textbf{98.49\%}
  & 0.493 & 0.494 & 0.481 \\
\textbf{RP（相对各自明文）}
  & 100\% & 98.49\% & ---
  & 100\% & \textbf{100.1\%} & \textbf{97.46\%} \\
\bottomrule
\end{tabular}
\end{center}

\vspace{0.6em}

\subsection*{WikiText-2 困惑度}

\begin{center}
\renewcommand{\arraystretch}{1.2}
\small
\begin{tabular}{@{}lccc@{}}
\toprule
\rowcolor{headerbg}
指标 & OSNIP 明文$\to$混淆 & fastProve structural & fastProve full \\
\midrule
PPL
  & $9.63\to12.58$
  & $11.59\to11.61$
  & $11.59\to11.73$ \\
相对增幅 $\Delta\%$
  & \textbf{$+31\%$}
  & \textbf{$+0.12\%$}
  & \textbf{$+1.15\%$} \\
\bottomrule
\end{tabular}
\end{center}

\vspace{0.7em}
\subsection*{结论}
\begin{enumerate}
  \item \textbf{任务保留率：} structural 在 10 项闭集上相对明文 RP $\approx 100\%$；
    full RP $\approx 97.5\%$，与 OSNIP 论文 $98.5\%$ 同量级。
  \item \textbf{PPL：} structural $+0.12\%$、full $+1.15\%$，
    远低于 OSNIP 报告的 $+31\%$（绝对 PPL 因窗口不同不可硬比）。
  \item \textbf{绝对分：} 本协议未使用 chat template，若干任务（PIQA / HellaSwag / MMLU）
    绝对准确率低于论文非私有基线；这不改变“混淆相对明文”的保留率结论。
  \item \textbf{未测：} 隐私 KNN-ASR、CNN/DM 摘要 ROUGE/BERTScore（OSNIP Table~3）。
\end{enumerate}

\vspace{0.3em}
{\scriptsize\color{note}
来源：OSNIP Tables~1--2；fastProve \texttt{osnip\_style\_bench\_n200.json}
（$n{=}200$/任务，BF16，RTX 5090，current covariant scheme）。
}
\end{document}
